\chapter{Physical System Setting}

\section{Architecture and participation semantics}
The network has $M$ distributed APs, each with $N_t$ transmit antennas, $K$
single-antenna UEs, and $P$ sensing targets. A central processor coordinates
beamforming and AP--target sensing participation over fronthaul. Communication
covariances $\{\mat{W}_k\}$ are globally cooperative. Dedicated sensing
covariances $\{\mat{S}_p\}$ are target-specific.

\begin{figure}[H]
\centering
\includegraphics[width=0.82\linewidth]{fig1_system_architecture.png}
\caption{Asymmetric architecture: globally cooperative communication and
target-specific dedicated sensing participation.}
\end{figure}

The binary variable $b_{mp}\in\{0,1\}$ authorizes AP $m$ to emit nonzero
dedicated sensing power for target $p$. It is therefore an AP--target dedicated
sensing-participation indicator. It does \emph{not} gate AP $m$'s communication
transmission. With a nonzero sensing-power floor, an authorized pair has a
physical sensing contribution; without that floor, $b_{mp}=1$ should be read
only as authorization rather than proof of nonzero emission.

\section{Signal and power model}
Let $\vect{h}_k$ be the stacked AP-to-UE channel and let
$\mat{W}_k=\vect{w}_k\vect{w}_k^H\succeq0$ denote the covariance of UE $k$'s
communication beamformer. The total transmit covariance is
\[
  \mat{R}_X=\sum_{k=1}^{K}\mat{W}_k+\sum_{p=1}^{P}\mat{S}_p.
\]
Because a UE does not in general know the dedicated sensing waveform,
$\sum_p\mat{S}_p$ enters its interference term unless waveform cancellation is
explicitly assumed. Each AP obeys
$\operatorname{tr}(\mat{E}_m\mat{R}_X)\leq P_{\max}$, whereas dedicated sensing
participation is enforced through
\[
 \operatorname{tr}(\mat{E}_m\mat{S}_p)\leq P_{\max}b_{mp}.
\]

\section{Default operating point}
The principal study uses a $400\,\mathrm{m}\times400\,\mathrm{m}$ area,
$M=6$, $N_t=2$, $K=3$, $P=2$, $N_\theta=2$, and $P_{\max}=0.1$ W (20 dBm).
The robustness radius is $\epsilon_h=0.05$ unless scanned. The cluster size is
varied as $N_{\rm req}\in\{2,3,4,5,6\}$.

\section{MISO versus massive MIMO}
At the UE level, the current downlink model is multi-user MISO: each UE has one
receive antenna, while the AP network provides a distributed multi-antenna
transmit array. ``Cell-Free massive MIMO'' describes the distributed-array
architecture and its scalable extension; it does not imply that the present
UE receiver is a multi-antenna MIMO receiver. The implementation deliberately
uses the MISO-equivalent observation model rather than a Kronecker MIMO
observation model requiring a different receiver architecture.

\sourcebox{Implementation sources: \texttt{sim/matlab/generate\_scenario.m},
\texttt{default\_params.m}, and \texttt{build\_E\_m.m}.}
